% ECONOMICS_PROBLEM: {"id": "C73-98", "title": "Pure positional equilibria in multiplayer safety games", "jel_code": "C73", "source_id": "OP-0162"}
% !TeX program = xelatex
\documentclass[11pt,a4paper]{article}
\usepackage[margin=23mm]{geometry}
\usepackage{fontspec}
\setmainfont{Latin Modern Roman}
\usepackage{amsmath,amssymb,mathtools}
\usepackage{xcolor,enumitem,booktabs,longtable,array,xurl,hyperref}
\definecolor{muted}{HTML}{53636C}
\hypersetup{colorlinks=true,urlcolor=blue,linkcolor=blue}
\setlength{\parindent}{0pt}
\setlength{\parskip}{5pt}
\newcommand{\R}{\mathbb R}
\newcommand{\E}{\mathbb E}
\newcommand{\N}{\mathbb N}
\newcommand{\ind}{\mathbf 1}
\newcommand{\Var}{\operatorname{Var}}
\newcommand{\Cov}{\operatorname{Cov}}
\newcommand{\supp}{\operatorname{supp}}
\DeclareMathOperator*{\argmin}{arg\,min}
\DeclareMathOperator*{\argmax}{arg\,max}
\newcommand{\entrymeta}[3]{{\sffamily\small\color{muted}\textbf{\texttt{#1}}\quad|\quad #2\par #3\par}}
\newcommand{\Problemref}[1]{\texttt{#1}}
\newcommand{\JELref}[2]{\texttt{#2} (JEL \texttt{#1})}
\begin{document}
\section*{Pure positional equilibria in multiplayer safety games}
\textbf{Problem ID:} \texttt{C73-98}\quad \textbf{JEL:} \texttt{C73}\par
% BEGIN ECONOMICS STATEMENT
\label{problem:OP-0162}
% GAME_THEORY_PROBLEM: {"local_id": "GT-032", "original_number": 32, "kind": "P", "entry_kind": "statement_only", "published": false, "review_required": true, "category": "game_theory"}
\label{game:GT-032}
{\sffamily\small\color{muted}
{\sffamily\small\color{muted}\textbf{GT-032} | Qualitative and graph games\par P | Source-reported pure positional existence question\par}
\par}
\subsection*{Definitions and mathematical statement}
A finite deterministic turn-based arena is a nonblocking directed graph $(V,E)$ with a partition $V=\bigsqcup_{i=1}^{n}V_i$; the owner of the current vertex chooses its successor. Player $i$ has a safe set $C_i\subseteq V$ and payoff
\[
 u_i^{v_0}(\sigma)=\Pr_{v_0,\sigma}(\forall t\geq0:\ v_t\in C_i).
\]
A pure positional profile consists of functions $f_i:V_i\to V$ selecting outgoing edges. The existence question is
\[
 \forall n\geq2\ \forall(V,E,(V_i),(C_i))\ \forall v_0\in V
 \quad\exists f\quad\forall i\quad\forall\tau_i,
 \qquad u_i^{v_0}(\tau_i,f_{-i})\leq u_i^{v_0}(f).
\]
Here $\tau_i$ may randomize and depend on the full finite vertex history.
\subsection*{Short English statement}
Can every finite multiplayer safety game have an equilibrium with deterministic choices based only on the current vertex?
\subsection*{Variants and scope boundaries}
For a history-dependent SPE, any player whose safe set was already left has continuation payoff zero, including after off-path histories. Requiring one positional NE from every fresh initial vertex is a different refinement. Safety is the complement of reachability for an individual event, but complementing every player's payoff reverses every player's preference and does not preserve nonzero-sum Nash inequalities. If all unsafe vertices are absorbing, the source gives an additional positive positional result.
\subsection*{Source relationship and status}
[S1], Section 5, separately retains safety-only pure positional existence. Its stationary randomized equilibrium results concern a larger strategy class. No stochastic transition vertices occur in the open pure positional formulation given here.
\subsection*{References}
\begingroup\footnotesize\raggedright
\textbf{[S1]} Mona Alluwaym, James C. A. Main and Sven Schewe, \emph{Simple Nash Equilibria for Qualitative Multiplayer Games}, MFCS 2026, LIPIcs 386, Article 85 (2026), doi:10.4230/LIPIcs.MFCS.2026.85. Locators: Section 2, pp. 85:4--85:5; Theorem 10 and Section 5, p. 85:14. \url{https://drops.dagstuhl.de/storage/00lipics/lipics-vol386-mfcs2026/LIPIcs.MFCS.2026.85/LIPIcs.MFCS.2026.85.pdf}
\par\endgroup

\subsection*{Common conventions: Source mathematical conventions}

\textbf{Finite games.} For a finite set $A$, $\Delta(A)$ is its probability
simplex. Players choose independent mixed strategies in Nash equilibrium;
correlation is allowed only when explicitly part of the solution concept.
In a game with payoff functions $u_i$ and strategy profile $x$, an additive
$\varepsilon$ Nash equilibrium satisfies
\[
 u_i(x)\geq u_i(x_i',x_{-i})-\varepsilon
 \quad\text{for every player $i$ and every admissible deviation $x_i'$}.
\]
For numerical approximation, payoffs are normalized as locally specified.
An $\varepsilon$ payoff residual is distinct from distance at most
$\varepsilon$ to the set of exact equilibria. Point convergence, convergence
of empirical play, and convergence in distance to an equilibrium set are
also distinct.

\textbf{Computational representations.} Unless stated otherwise, a complexity
question takes rational data with binary encoding; polynomial time is measured
in total bit length. A fixed-accuracy polynomial algorithm is not necessarily
polynomial in $1/\varepsilon$, and neither is necessarily polynomial in
$\log(1/\varepsilon)$. Succinct games and value, demand, or sampling oracles
must declare their interfaces. Exact real solutions require an explicit
representation; an unspecified list of arbitrary real numbers is not a
finite computational output. A claimed algorithmic barrier is meaningful
only for its specified model and reductions.

\textbf{Dynamic games.} Histories, information, observation, and allowable
randomization are part of the model. A stationary strategy depends only on
the current state; a positional strategy in a deterministic graph game
chooses one action at each controlled vertex. Nash equilibrium at an initial
state and equilibrium after every history are different requirements.
An expectation of a pathwise $\liminf$ payoff, a $\liminf$ of expected averages,
a limiting expected average, and a uniform finite-horizon equilibrium are
different criteria. None is silently substituted for another.

\textbf{Allocation and mechanisms.} A complete allocation assigns every item
exactly once unless otherwise stated. Zero-valued goods, negative values,
capacity constraints, feasibility, lotteries, and copies of goods affect
fairness definitions and are specified locally. Dominant-strategy incentive
compatibility, Bayesian incentive compatibility, universal truthfulness,
and truthfulness in expectation impose different inequalities. Whenever
an objective ratio has zero denominator, its convention is stated or those
instances are excluded.

\textbf{Infinite-dimensional models.} Expectations, integrals, stochastic
equations, and optimization objectives require their local regularity,
integrability, filtrations, and admissible domains. A backward--forward
mean-field system is a boundary-value problem; it is not an initial-value
semiflow without an additional construction. Long-time, large-population,
and vanishing-noise limits require an order of limits and a convergence
topology. If a minimum need not be attained, the objective is its infimum
or supremum and the approximation requirement is stated separately.
% END ECONOMICS STATEMENT
\end{document}
